\section{Background}
\label{sec:bg}

This section fixes the vocabulary used by the rest of the paper. A reader
fluent in RoCEv2 and \texttt{mlx5} can skip to \S\ref{sec:measurement}.

\subsection{RoCEv2 transport}

RDMA over Converged Ethernet version~2 carries an InfiniBand transport packet
inside a UDP datagram, over IP, over Ethernet~\cite{ibta-a17}. The UDP
destination port is fixed at 4791 and the source port is used as an
entropy field for equal-cost multipath hashing, so a queue pair is identified
on the wire by a stable five-tuple.

The unit of communication is the \emph{queue pair} (QP): a send queue and a
receive queue whose contexts on the two endpoints are bound to each other. A
reliable connected (RC) queue pair guarantees ordered, lossless delivery to the
consumer; an unreliable datagram (UD) queue pair does not. Each RC packet
carries a \emph{packet sequence number} (PSN) in its base transport header
(BTH). The responder checks the PSN of every arriving packet against the one it
expects: a match advances the expected PSN and, for a write or a send, places
the payload; a gap causes the responder to emit a negative acknowledgement
(NAK) naming the first missing PSN and to discard everything until the
requester restarts from there. Recovery is therefore \emph{go-back-N}: one lost
packet costs the retransmission of every packet after it in the current
message, which is why the goodput cost of loss scales with message size rather
than with the number of lost packets. The requester also runs a local
acknowledgement timeout; if no acknowledgement arrives within it, the requester
retransmits from the last acknowledged PSN. Extended headers carry the
operation-specific fields: RETH for the remote address and length of an RDMA
write or read request, AETH for the acknowledgement syndrome and credit count,
DETH for the source queue pair of a datagram. An invariant cyclic redundancy
check (ICRC) covers the packet with the mutable fields masked out.

Two facts from this structure matter throughout. First, an RC requester is
allowed to have many packets outstanding, and the number it actually keeps
outstanding is a device property that the specification does not fix. Second,
the only signal a sender gets about receiver-side loss is the absence of an
acknowledgement or the arrival of a NAK, so a receiver that discards below the
transport layer, at its physical-layer ingress, is invisible to the sender
except through the retransmission it causes.

\subsection{The \texttt{mlx5} software interface}

The device is programmed through rings in host memory and a small number of
memory-mapped doorbell pages~\cite{mlx-prm,rdma-core}.

A \emph{send queue} (SQ) and \emph{receive queue} (RQ) are contiguous rings of
work-queue-element (WQE) slots in host memory, sized in units of a basic block.
The consumer posts a WQE describing an operation (opcode, scatter-gather list,
remote address, flags, whether it is signalled) and advances a producer index.
A \emph{completion queue} (CQ) is a third ring, written by the device; each
completion-queue entry (CQE) carries an ownership bit that flips on every wrap,
so the consumer polls for ownership rather than for an interrupt. Unsignalled
work requests produce no CQE at all, and their slots are reclaimed by a later
signalled completion.

Publication has two steps. The provider first writes the new producer index
into a \emph{doorbell record}, an eight-byte field in host memory that the
device reads by direct memory access. It then writes a doorbell into a
\emph{user access region} (UAR), a page of device memory mapped into the
process, using a write-combining store. Small work requests may be copied
entirely into the UAR page instead, so the device never has to fetch them; the
provider calls this the BlueFlame path. Because a single doorbell can publish
an arbitrary prefix of newly posted work requests, the number of memory-mapped
writes per work request falls as the posting batch grows, and the device sees
work in batches rather than one at a time.

The device holds per-queue-pair state in a \emph{queue-pair context} (QPC):
transport type, state, path attributes including the destination global
identifier and the path maximum transmission unit, the send and receive PSNs,
retry and timeout settings, and the congestion-control state. Memory
registration is described by a \emph{memory protection table} (MPT) entry
naming a region, its keys and its access rights, and a \emph{memory
translation table} (MTT) giving the virtual-to-physical mapping. Both are
structures in host memory that the device caches on chip.

The queue-pair state machine runs RESET to INIT to RTR (ready to receive) to
RTS (ready to send), driven either by a manual exchange of queue-pair number,
PSN, global identifier and path attributes over a side channel, or by the
connection manager. Errors move the queue pair to the error state, after which
it must be reset.

\subsection{DCQCN}

Data Center Quantized Congestion Notification~\cite{dcqcn2015} is the
rate-based congestion-control algorithm this device family implements in
firmware. It has three parts.

A \emph{congestion point} marks packets. In the textbook deployment this is a
switch running random early detection on its egress queue: a packet that is
ECN-capable, that is, one whose IP header carries ECT(0) or ECT(1), is
rewritten to congestion-experienced (CE) with a probability that rises from
zero at a low threshold $K_{\min}$ to $P_{\max}$ at a high threshold
$K_{\max}$.

A \emph{notification point} (NP) at the receiver turns CE-marked data packets
into congestion notification packets (CNPs) sent back to the source queue pair,
rate-limited to at most one per configured minimum interval per queue pair.

A \emph{reaction point} (RP) at the sender keeps per-queue-pair state: a
current rate $R_C$, a target rate $R_T$ and a congestion parameter $\alpha$. On
a CNP the reaction point updates $\alpha$ toward one, cuts $R_C$ by a factor
involving $\alpha$, and remembers the pre-cut rate as the target. In the
absence of CNPs it recovers, first quickly toward the target (fast recovery),
then by additive increase, then by hyper-increase, driven by two clocks: a
timer and a byte counter, both of which must fire before a stage advances.

Two properties matter for what follows. The reaction point's state is
\emph{persistent per queue pair} and survives across work requests, so the rate
a message gets depends on what the previous messages on that queue pair did.
And the whole loop is invisible in the data path: the only endpoint-visible
evidence that it is running is the counter pair \ctr{np\_cnp\_sent} at the
receiver and \ctr{rp\_cnp\_handled} at the sender.
